\documentclass[12pt]{article}
\usepackage{amssymb,amsmath, endnotes,setspace,enumerate,ulem,color,xspace,lscape}
\usepackage{amsthm,subfigure,graphics,epsfig, hyperref}
\usepackage[sort]{natbib}
\allowdisplaybreaks[1]
\normalem
\newcommand{\textR}[1]{\textcolor{blue}{\texttt{#1}}}
\newcommand{\R}{\textR{R}}

\begin{document}

\title{Title of project}  

\author{
Adam Field   \thanks{Department of Physics.  Email: {\tt adfield@wpi.edu}}  
}

\maketitle

\begin{center}
\textbf{Abstract}
\end{center}

This is the abstract of the report \dots

\thispagestyle{empty}

\newpage

\pagenumbering{arabic} % add page numbers to the report

\section{Introduction}\label{intro}

Need to provide a roadmap at the end of Section \ref{intro}. 
 
\section{Methodology}

Do not repeat what has been covered in the lecture notes. 

Use this example to see how to cite a paper: \cite{rizzo2007statistical}

\subsection{Subsection 1}

\subsection{Subsection \dots}

\section{Simulation Study}	

\subsection{Subsection 1}

\subsection{Subsection \dots}


\section{Real Data Analysis}


\section{Conclusions}

 
\newpage


% Reference
\bibliographystyle{apalike}
\bibliography{referencelist}


\newpage

% Appendix
\appendix

\begin{center}
{\Large {\bf Appendix: \R\ code}}
\end{center}

% INSERT R CODE HERE, 10 point font
{\footnotesize \begin{verbatim}
m <- 10000 
theta.hat <- se <- numeric(5)
g <- function(x) {
  exp(-x - log(1+x^2)) * (x > 0) * (x < 1)
}

x <- runif(m)     #using f0
fg <- g(x)
theta.hat[1] <- mean(fg)
se[1] <- sd(fg)
\end{verbatim} }



\end{document}


